\section{البنية العامة لنظام التحكم الصوتي}
لضمان استجابة ديناميكية دقيقة وتجنب أي تداخل مع حلقات التحكم الأساسية، تم تصميم نظام التفاعل الصوتي ليعمل كنظام مدمج (\textenglish{Embedded System}) مستقل في معالجة الإشارة الصوتية وتفسيرها، على أن يكتفي بإرسال الأوامر التنفيذية النهائية الصافية إلى المتحكم الرئيسي للذراع الروبوتية.

تتكون السلسلة الوظيفية للهيكلية الهرمية المصممة لنظام التحكم الصوتي من المراحل المتتالية التالية:
\begin{enumerate}
    \item \textbf{التقاط الإشارة الصوتية:} يتم استخدام وحدة التقاط صوتية مخصصة لتحويل الموجات الصوتية الميكانيكية الناتجة عن الكلام البشري إلى إشارات كهربائية مستمرة، والتي تخضع لاحقاً لعمليات الترشيح والرقمة لتصبح جاهزة للمعالجة.
    \item \textbf{وحدة المعالجة المركزية (التحليل الصوتي):} تستقبل الإشارة الرقمية الممثّلة للأمر الصوتي وتمررها عبر خوارزميات التعرف على الكلام وتحويل الصوت إلى نص (\textenglish{Speech-to-Text}) لاستخلاص الكلمات والعبارات المطابقة للنماذج اللغوية والصوتية المدمجة.
    \item \textbf{مفسر الأوامر (\textenglish{Command Mapper}):} خوارزمية برمجية تقوم بالبحث عن الكلمات المفتاحية (\textenglish{Keywords}) المحددة ضمن النص المستخرج وتصفيتها من الكلمات اللغوية الزائدة، ومن ثم ربط هذه الكلمات بمتجهات حركية أو أوامر تشغيلية محددة.
    \item \textbf{بروتوكول الاتصال (\textenglish{Communication Link}):} يتم نقل الأمر الرقمي المفسّر (مثل تعديل إحداثيات الهدف الكارتزي أو تغيير حالة الماسك) عبر قناة اتصال تسلسلية موثوقة ومحمية من الضجيج إلى معالج الذراع الرئيسي.
    \item \textbf{تنفيذ الحركة الميكانيكية:} يستقبل معالج الذراع الأمر ضمن طبقة التحكم عالي المستوى (\textenglish{HLC})، حيث يتم تحديث شعاع الموضع المرجعي ($X_{\text{ref}}$) وحساب الزوايا المفصلية المرجعية ($q_{\text{ref}}$) عبر بلوك الحركيات العكسية (\textenglish{IK})، لتتولى بعد ذلك طبقة التحكم منخفض المستوى (\textenglish{LLC}) توليد إشارات القيادة ومزامنة محركات السيرفو الخمسة لتحقيق الحركة المطلوبة.
\end{enumerate}